\section{Las Palabras del Lenguaje - Variables y Tipos de Datos}
\label{sec:variables_tipos}

En el Capítulo 1, dimos nuestros primeros pasos en el mundo de C++ con el programa "Hola, Mundo". Fue una introducción a cómo un programa puede interactuar con la consola. Pero los programas no solo muestran mensajes; su verdadero poder reside en la capacidad de procesar y manipular información. Para ello, necesitamos un lugar donde guardar esos datos.

\subsection{Qué es una Variable: La Caja con Etiqueta}
\label{subsec:que_es_variable}

Imagina que estás organizando una mudanza. Tienes muchos objetos diferentes (libros, ropa, utensilios de cocina) y necesitas guardarlos en cajas. Para no olvidar qué hay en cada caja, le pones una etiqueta: "Libros de cocina", "Ropa de invierno", etc.

En programación, una \textbf{variable} es exactamente como una de esas cajas con etiqueta. Es un espacio en la memoria del ordenador donde podemos guardar un dato. La "etiqueta" de la caja es el \textbf{nombre de la variable}, y el "contenido" de la caja es el \textbf{valor} que guarda.

Por ejemplo, si queremos guardar la edad de una persona, podemos crear una variable llamada `edad` y asignarle el valor `30`. Si luego la persona cumple años, podemos cambiar el valor de esa variable a `31`.

Las variables son fundamentales porque nos permiten:
\begin{itemize}
    \item \textbf{Almacenar datos:} Desde números y texto hasta estructuras más complejas.
    \item \textbf{Manipular datos:} Podemos cambiar el valor de una variable a lo largo de la ejecución del programa.
    \item \textbf{Hacer nuestros programas dinámicos:} En lugar de trabajar con valores fijos, podemos procesar información que cambia (como la entrada del usuario o los resultados de cálculos).
\end{itemize}

En las siguientes secciones, exploraremos los diferentes tipos de "cajas" que C++ nos ofrece, es decir, los tipos de datos que podemos guardar en nuestras variables.
